\section{Entropy and counting consequences}
\label{cons:section}

The pointwise entropy bound has two kinds of counting consequences.
Maximizing a separable entropy objective estimates the total matching
count, while evaluating it at a feasible uniform mean gives lower counts
from structural hypotheses. Throughout this section, unless
an empty graph is explicitly mentioned, \(G\) is a finite loopless
multigraph on \(2m>0\) vertices with at least one perfect matching.
Parallel edges retain their individual labels. Write \(P=P(G)\),
\(N=|\mathcal M(G)|\), and recall that \(F\), \(B\), and \(H\) use
natural logarithms. For every \(x\in P\), Theorem~\ref{ent:main} gives
\begin{equation}
 \label{cons:input}
 F(x)-\left(2-\frac2m\right)B(x)\le H(x)\le F(x).
\end{equation}
The elementary scalar estimate
\begin{equation}
 \label{cons:complement}
 0\le -(1-u)\ln(1-u)\le u\qquad(0\le u\le1)
\end{equation}
will be used repeatedly. Indeed,
\(u+(1-u)\ln(1-u)\) vanishes at zero and has derivative
\(-\ln(1-u)\ge0\) on \([0,1)\); continuity gives the endpoint.
Every perfect matching has \(m\) edges, so \(\sum_e x_e=m\) and
\begin{equation}
 \label{cons:B-bound}
                         0\le B(x)\le m.
\end{equation}

\subsection{Entropy defects and separable counting estimates}

\begin{corollary}[Entropy-dependent defect]
\label{cons:defect}
For every \(x\in P\),
\begin{equation}
 \label{cons:defect-bound}
 0\le F(x)-H(x)
 \le \min\{F(x),\,2m-2\}.
\end{equation}
\end{corollary}

\begin{proof}
The first inequality follows from \(H\le F\). Nonnegativity of
entropy gives \(F-H\le F\), while
Equations~\eqref{cons:input} and~\eqref{cons:B-bound} give
\(F-H\le(2-2/m)B\le2m-2\).
This includes \(m=1\), when the defect is zero.
\end{proof}

Define the binary-coordinate entropy by
\[
 H_{\mathrm{bin}}(x)
   =\sum_e[-x_e\ln x_e-(1-x_e)\ln(1-x_e)]
   =F(x)+B(x).
\]
The word ``binary'' refers to an edge indicator's two values;
these entropies are still measured with natural logarithms.

\begin{corollary}[Separable estimates of the matching count]
\label{cons:entropy-count}
One has
\begin{align}
 \ln N&\le\max_{x\in P}F(x)\le\ln N+2m-2,
       \label{cons:F-count}\\
 \ln N&\le\max_{x\in P}H_{\mathrm{bin}}(x)
                         \le\ln N+3m-2.
       \label{cons:binary-count}
\end{align}
\end{corollary}

\begin{proof}
The entropy comparison in Equation~\eqref{ent:comparison}, applied to
the uniform law, gives entropy at most \(\ln N\) for every law on
\(N\) outcomes, with equality for the uniform law.
At the mean \(\bar x\) of the uniform matching law we therefore have
\(H(\bar x)=\ln N\). Equation~\eqref{cons:input} gives
\(\ln N\le F(\bar x)\le H_{\mathrm{bin}}(\bar x)\), proving both
lower estimates. Corollary~\ref{cons:defect} gives
\(F(x)\le H(x)+2m-2\le\ln N+2m-2\) at every \(x\in P\).
Also,
\[
 H_{\mathrm{bin}}(x)-H(x)
 \le\left(3-\frac2m\right)B(x)\le3m-2.
\]
Taking maxima proves the upper estimates. The maxima exist because
\(P\) is compact and both separable functions are continuous.
\end{proof}

\subsection{Uniform marginals and structural lower counts}

\begin{theorem}[Many edge-disjoint perfect matchings]
\label{cons:structural}
Let \(G\) be a finite loopless multigraph on \(2m>0\) vertices that
contains \(k\ge1\) pairwise edge-disjoint perfect matchings.
If \(k=1\), then \(N(G)\ge1\). If \(k\ge2\), then
\begin{equation}
 \label{cons:finite-order-count}
 N(G)\ge k^m(1-1/k)^{2(m-1)(k-1)}
       \ge e^{-2(m-1)}k^m
       \ge(e^{-2}k)^m.
\end{equation}
\end{theorem}

\begin{proof}
Let \(M_1,\ldots,M_k\) be the selected matchings and take their
uniform mixture. Its mean
\[
                 x=\frac1k\sum_{i=1}^k\mathbf1_{M_i}\in P(G)
\]
is \(1/k\) on their union of \(km\) labeled edges and zero on every
other edge. In particular, the hypotheses do not require all ambient
edges to belong to that union. We have \(F(x)=m\ln k\).
If \(k=1\), the selected matching already gives \(N(G)\ge1\).
For \(k\ge2\), direct evaluation yields
\[
                 B(x)=-m(k-1)\ln(1-1/k).
\]
Since \(H(x)\le\ln N(G)\), Theorem~\ref{ent:main} now gives
\[
 \ln N(G)\ge m\ln k+2(m-1)(k-1)\ln(1-1/k).
\]
Exponentiating proves the first inequality in
Equation~\eqref{cons:finite-order-count}.
The function \(\ln(1-u)+u/(1-u)\) starts at zero and has derivative
\(u/(1-u)^2\ge0\) for \(0\le u<1\). Thus
\((k-1)\ln(1-1/k)\ge-1\), giving the second inequality; the third
uses \(m-1\le m\).
\end{proof}

At \(m=1\), the finite-order estimate is simply
\(N(G)\ge k\), the exact count on the selected parallel-edge subgraph.
We keep \(k=1\) separate to avoid an undefined power at zero.

Only feasibility of the uniform mean enters the preceding proof. It
therefore applies to the following class without requiring a partition
of its edges into perfect matchings.

\begin{corollary}[Regular graphs with large odd cuts]
\label{cons:regular}
Let \(G\) be a finite loopless \(k\)-regular multigraph on \(2m>0\)
vertices, where \(k\ge1\). Suppose
\(\lvert\delta(S)\rvert\ge k\) for every odd vertex set \(S\).
Then the conclusions of Theorem~\ref{cons:structural} hold.
In particular, they hold for every \((k-1)\)-edge-connected
\(k\)-regular graph of positive even order.
\end{corollary}

\begin{proof}
The vector \(x_e=1/k\) has unit degree sums, and the cut hypothesis
gives \(x(\delta(S))\ge1\) for every odd \(S\).
Theorem~\ref{geom:thm-polytope} implies \(x\in P(G)\); in particular,
\(G\) has a perfect matching. Since \(|E|=km\), the evaluations of
\(F(x)\) and \(B(x)\) in the preceding proof apply unchanged when
\(k\ge2\). For \(k=1\), the graph itself is a perfect matching.

For the final assertion, an odd \(S\) is nonempty and proper because
\(|V(G)|\) is even. Edge connectivity gives
\(\lvert\delta(S)\rvert\ge k-1\), while summing degrees in \(S\)
gives
\[
            k|S|=2|E(S)|+|\delta(S)|.
\]
Consequently \(|\delta(S)|\equiv k\pmod2\); it cannot equal \(k-1\)
and is therefore at least \(k\). Parallel edges are counted with
multiplicity in this identity.
\end{proof}

The odd-cut hypothesis is a condition on the uniform vector. Requiring
merely that every edge lie in some perfect matching does not assert
that condition. At small degree the exponential base in
Equation~\eqref{cons:finite-order-count} can be below one, so these
estimates do not supply the classical cubic exponential lower bound.

The empty graph has one empty perfect matching. Its count is handled
directly, without using expressions that divide by \(m\); for
\(k\ge1\), the homogeneous coarse bound \((e^{-2}k)^m\)
is interpreted as one.
When \(k=0\) and the vertex set is nonempty, the coarse
bound is zero, while formulas containing \(1/k\) are not asserted.
In particular, specifying no selected matchings places no restriction
on the ambient edge set.

The weaker separable and nonlinear comparisons and their structural
consequences appear with the alternative proofs in Appendix~\ref{a:section}.
The next section examines sharp face codimension and the obstruction to
coefficient one. Section~\ref{cons:weighted-section} then passes from
these entropy estimates to weighted partition functions.
